\section{Introduction: From Hand-Crafted Sketches to Learned Sufficiency}\label{sec:min-introduction}

Long-document measurement is usually presented as a scoring problem: given a
manifesto, report where a party stands on economic policy, social liberalism,
immigration, European integration, environmental policy, or decentralization.
For a language-model system, however, the more basic object is a compression
problem. A long document is split, shortened, retrieved from, summarized, or
otherwise transformed before a final readout is applied. The final score is
therefore not only a property of the original document. It is also a property
of the representation that survived the path from raw text to the model's
context window.

That distinction matters for political methodology. Fiscal commitments, tax
qualifications, public-service promises, market-intervention claims, and
welfare retrenchment cues often appear in different parts of a manifesto. The
document-level judgment depends on their joint interpretation. A compression
step can lose the caveat that changes a tax pledge, overweight a rhetorically
salient paragraph, or erase the boundary fact that turns two neighboring
sections into a coherent policy position. If the final model reads the
compressed object faithfully, it can still answer the wrong question because
the compressed object is not sufficient for the full-document oracle.

The core question is therefore:
\[
  \text{when is the state that reaches the root sufficient for the target
  readout?}
\]
C-TreePO answers this question by treating recursive LLM summarization as a
learned mergeable sketch. The word \emph{sketch} is doing real work here, not
only supplying an analogy. A sketch is a small state that is designed so a
query can be answered after local states are merged. The state must contain
the information the query needs, and the merge path must preserve that
information.

The simplest example is Count-Min \citep{CormodeMuthukrishnan2005}. Given a
stream or multiset $A$, a hand-designed state map $g$ stores a small
$d\times w$ matrix of counters. Each item increments one counter per row.
Two disjoint stream segments merge by adding their counter matrices:
\[
  g(A\uplus B)=g(A)+g(B).
\]
For an item $q$, the readout $f_q$ estimates its frequency by checking the
counters to which $q$ hashes and taking the row minimum:
\[
  f_q(g(A))=\min_{j\le d} g(A)_{j,h_j(q)}.
\]
The merge is exact at the state level. The readout is approximate only because
other items can collide with $q$ in the same counters. This is the classical
template: choose a fixed query family, hand-design a state function $g$ and an
associative merge, then prove that the readout remains valid after parallel or
tree-shaped aggregation. HyperLogLog, quantile summaries, Bloom filters, and
many streaming summaries follow the same pattern
\citep{FlajoletEtAl2007,Agarwal2013MergeableTODS,BoykinEtAl2014Summingbird}.

This template is also a sufficiency statement. In Fisher's language, the state
is a statistic that preserves all information relevant to the target
\citep{Fisher1922}; in Blackwell's language, the compressed experiment is
equivalent to the original for the decision problem induced by the query
\citep{Blackwell1953}. The Count-Min counter matrix is not a lossless copy of
the stream. It deliberately discards order, syntax, and identities not needed
for the chosen frequency query. Its success comes from being sufficient for
that query up to a controlled collision residual.

Classical sketches get this property by hand-crafting both sides of the
interface. The readout $f$ is chosen from a small catalog: frequency,
cardinality, quantile, membership, or a related closed-form query. The state
function $g$ is engineered to match that query: counters for frequencies,
register maxima for cardinality, ordered tuples for quantiles, bit arrays for
membership. The proof is algebraic because the designer knows what state is
sufficient before the algorithm runs.

C-TreePO keeps the interface but changes where the design burden sits. In
long-document measurement, the target readout is not a closed-form streaming
query. It may be a trusted human coding rule, a calibrated rubric judge, a
downstream loss, or an expert-survey target. We write that target-facing
readout as $f$. The sufficient state is not known in advance. C-TreePO learns
a state map $g$ that tries to make
\[
  f(g(x)) \approx f(x)
\]
on the objects the tree actually creates. For a binary merge, the local
identity becomes
\[
  f\!\left(g(g(u)\concat g(v))\right)
  \approx
  f(u\concat v).
\]
This is the mergeable-sketch identity in semantic form: summarize the left
span, summarize the right span, merge those states, and recover the same
readout that would have been obtained from the raw union.

The move from hand-crafted sketches to C-TreePO is a move from analytic state
design to learned state design. In an LLM application, $g$ may be a prompted
summarizer, a DSPy/GEPA-optimized prompt program, an instruction-tuned model,
or a weight-tuned model trained against oracle labels and local-law feedback.
The readout $f$ may likewise be a prompted scorer, a calibrated judge, or a
learned evaluator. In a non-language domain, the same pair may be realized by
ordinary neural networks, transformers over structured states, or neural
operators such as Fourier neural operators that learn mappings between
function spaces \citep{Kovachki23NeuralOperator,Li21FNO}. The architecture is
not the definition. The definition is the state--merge--readout contract and
the empirical certificate that the contract held on the deployed tree.

This is also where LLMs make the learned-sketch problem unusually attractive.
Classical sketch states are often opaque to domain experts. A Count-Min matrix
can be audited for dimensions, hashes, and collision bounds, but a political
scientist cannot read a counter matrix and see which fiscal qualification it
contains. An LLM state $g(x)$ is natural language. It can name the economic
pledges, quote or paraphrase qualifications, record uncertainty, and expose
what evidence the downstream readout will see. The summary can be inspected,
compared with the source span, and stress-tested by counterfactual examples
that share the same summary but should or should not share the same oracle
value. Observability does not prove sufficiency, but it makes sufficiency
auditable in a way numeric sketches are not.

The audit is necessary because LLM pipelines otherwise create black-box
compression. A long-context model may score a whole document directly, but it
can still use information unevenly, including the familiar
``lost-in-the-middle'' pattern \citep{LiuEtAl2024LostMiddle}. A retrieval or
recursive-summary pipeline is more transparent operationally, but it can hide
which evidence was dropped and which cross-span interactions survived. Either
way, the representation being scored must be checked against the oracle, not
only trusted because the final prediction is plausible.

Design-Based Supervised Learning supplies the measurement discipline for that
check. DSL's warning is that model-produced labels can have high accuracy and
still induce biased downstream estimands when errors align with the target
analysis \citep{EgamiEtAl2024}. C-TreePO applies the same logic one level
earlier, to the representation being labeled. Root labels anchor the
full-document target. Sampled leaf and merge labels estimate whether the
compression path preserved the target. Logged inclusion probabilities make the
sampled checks a design-based estimate of the finite population of local
failures rather than a convenience sample of examples the researcher happened
to inspect.

The method therefore has two linked pieces. The first is algebraic: split a
document into contiguous leaves, map each leaf to a state with $g$, recursively
merge adjacent states with $g$, and apply $f$ at the root. Three local laws
ask whether each step preserved the task-relevant state. Leaf summaries must
be sufficient for their source spans, re-summarizing stored states must not
drift, and merging child states must preserve the same cross-span information
as summarizing the raw union. Under exact laws, the root state is sufficient
for the full-document readout. Under approximate laws, the accumulated local
residual supplies a certificate for compression-induced distortion.

The second piece is design-based. A realized tree contains roots, leaves,
stored summaries, and internal merge nodes. These are finite-population units.
The researcher can sample them with logged positive propensities, query a
gold-standard coder or high-grade judge on the sampled checks, and use
IPW/Horvitz--Thompson estimators to estimate tree-wide violation rates and
distortion. This changes the supervision problem. Instead of spending every
expensive label on a full document, the researcher can allocate some expert
attention to cheaper partial-document units, as long as the local-law audit
certifies that those units measure the same target-preservation property.

The Manifesto/Benoit application is the running example. The target is a set
of document-level policy positions. The representation is a tree of
policy-evidence summaries. The readout $f$ is a rubric-facing scorer for a
policy dimension. The state map $g$ is an LLM summarizer trained or prompted
to preserve the evidence that the rubric uses. On the six-dimension benchmark,
a per-dimension C-Tree pipeline using a single open-weight model reaches macro
Pearson $r=0.829$ against expert-survey means at $8\mathrm{K}$-character
leaves, above Benoit's proprietary 18-score ensemble ($0.817$) and their
matched Gemma-3 open-weight replication ($0.793$)
\citep{BenoitEtAl2025,ManifestoProject2025a}. The economic dimension is the
detailed granularity case: the $8\mathrm{K}$-character tree reaches
$r=0.939$ against expert means, and an alternating prompt ladder shows a
narrow plateau across $1024$--$8192$ token leaves.

The empirical results are important, but the paper's methodological claim is
the frame: LLM compression can be treated as a learned, observable,
mergeable-sketch problem. Classical sketches verify sufficiency by proof
because the state is hand-designed. C-TreePO verifies approximate sufficiency
by local law, design-based audit, and direct inspection of the learned state.
The rest of the manuscript builds outward from that claim.
Section~\ref{sec:min-framework} gives the state interface, Count-Min warm-up,
compression tree, and local laws. Section~\ref{sec:min-audit} turns those laws
into a partial-document supervision design. Section~\ref{sec:min-guarantees}
states the conditional theorem stack. Section~\ref{sec:min-manifesto-results}
reports the Manifesto results. Section~\ref{sec:min-validation} connects the
application back to controlled sketch, Markov, and neural-operator anchors.
Section~\ref{sec:min-applications} closes with preference training,
distillation, and scope.
